\scanpage{14}\mspage{13}
$=$ comp. connexe que les $\{\{x\}\}$, comme on voit à l'aide de le multistrate
\[
\bigl\{\{\{x\}\},\{L\setminus\{x\}\}\bigr\}
\]

Si $L$ \deleted{de card. $1$,} $=\emptyset$, \deleted{$\unsure{\mathcal P|\mathfrak P}(L)$ est de} \deleted{cardi} alors $M=\emptyset$, et $M$ est connexe.

Si $\operatorname{Card}L=1$, $\mathcal P(L)=\{\emptyset,L\}$, et les $\unsure{\mathscr{L}|\Phi}\subset\mathcal P(L)$ \deleted{qui sont \unsure{des} sont que} \inserted{telles} \deleted{\unsure{$\{\mathscr{L}\}$} qui ne contiennent pas $\emptyset$} pour $\emptyset\notin\mathscr{L}$ sont $\emptyset$ et $\{L\}$. Donc $M=\{\{L\}\}$ a un seul élément, et est connexe.

Donc

\uline{Prop} Si $\operatorname{card}L\leq1$, $M$ est connexe (il est vide si $\deleted{\unsure{\mathrm{Card}}}L=\emptyset$, il est de card $1$ donc $0$-connexe si $\operatorname{card}L=1$). Si $\operatorname{card}L\geq1$, alors les $\deleted{\unsure{Y|F}}X\in M$ qui ont un élément minimal sont tous dans la $=$ comp. connexe. Donc si on se borne aux \deleted{familles} figures ensemblistes \uline{finies} (ce qui n'offre pas d'inconvénient à mes yeux) $M$ est connexe. Donc l'axiome C4 \deleted{est ici de} relatif à la décomposition de $M$ en comp. connexes (que $X,Y$ appartenant à des composantes différentes sont compatible) est automatique.
